\subsection{Data}

The empirical foundation of this study is a dataset of 424 United States-listed equities and exchange-traded funds spanning ten years (2014--2024).
We focus our single-asset analysis on the SPDR S\&P~500 ETF Trust (SPY), which tracks the S\&P~500 index.
The data include daily open, high, low, and close prices along with volume-weighted average price (VWAP) metrics.
The in-sample dataset comprises $T = 2{,}766$ trading days (January~3, 2014 to December~31, 2024).
A subsequent 219 trading days in 2025 are reserved for out-of-sample testing.

\subsection{Excess Growth Rate Calculation}

We compute the excess growth rate $G_{i,j}$ for ticker $i$ between time periods $j-1 \to j$ given the equity price series $P_{i,j}$:
\begin{equation}
    G_{i,j} \equiv \frac{1}{\Delta t} \cdot \ln\left(\frac{P_{i,j}}{P_{i,j-1}}\right) - r_f
    \label{eq:growth_rate}
\end{equation}
where $\Delta t = 1/252$ (one trading day in years) and $r_f$ is a constant risk-free rate.
The excess growth rate has units of year$^{-1}$ and is time-additive under continuous compounding.

\subsection{Continuous Hidden Markov Model}

We define the hidden Markov model as the tuple $\mathcal{M} = (\mathcal{S}, \mathcal{O}, \mathbf{T}, \mathbf{E}, \bar{\pi})$ \citep{rabiner1986introduction}, where:
\begin{itemize}
    \item $S_t \in \mathcal{S} = \{1, \ldots, K\}$ are the hidden states representing market regimes,
    \item $\mathcal{O} \subset \mathbb{R}$ is the continuous observation space of excess growth rates,
    \item $\mathbf{T} \in \mathbb{R}^{K \times K}$ is the state transition matrix with $T_{ij} = \Prob(S_{t+1} = j \mid S_t = i)$,
    \item $\mathbf{E} = \{(\mu_k, \sigma_k)\}_{k=1}^{K}$ specifies the Gaussian emission distributions $G_t \mid S_t = k \sim \mathcal{N}(\mu_k, \sigma_k^2)$,
    \item $\bar{\pi}$ is the stationary distribution satisfying $\bar{\pi} = \bar{\pi}\mathbf{T}$.
\end{itemize}

Unlike the discrete approach of \citet{alswaidan2026hybrid}, which discretizes observations into Laplace quantile bins and estimates $\mathbf{T}$ by frequency counting, we learn $\mathbf{T}$, $\{\mu_k\}$, $\{\sigma_k\}$, and $\bar{\pi}$ jointly from the continuous observations via the Baum-Welch algorithm.

\subsection{Baum-Welch (EM) Algorithm}
\label{sec:baum_welch}

The Baum-Welch algorithm \citep{baum1970maximization} iterates between two steps until convergence.

\paragraph{E-Step.}
Compute the forward variables $\alpha_t(k) = \Prob(O_1, \ldots, O_t, S_t = k \mid \mathcal{M})$ and backward variables $\beta_t(k) = \Prob(O_{t+1}, \ldots, O_T \mid S_t = k, \mathcal{M})$ via:
\begin{align}
    \log \alpha_t(j) &= \text{logsumexp}\left(\log \alpha_{t-1}(\cdot) + \log T_{\cdot,j}\right) + \log b_j(O_t) \label{eq:forward} \\
    \log \beta_t(i) &= \text{logsumexp}\left(\log T_{i,\cdot} + \log b_{\cdot}(O_{t+1}) + \log \beta_{t+1}(\cdot)\right) \label{eq:backward}
\end{align}
where $b_k(O_t) = \phi(O_t; \mu_k, \sigma_k)$ is the Gaussian emission density.
All computations are performed in log-space to prevent numerical underflow for long sequences \citep{bilmes1998gentle}.

The posterior state probability is:
\begin{equation}
    \gamma_t(k) = \Prob(S_t = k \mid O_{1:T}, \mathcal{M}) = \frac{\alpha_t(k) \beta_t(k)}{\sum_{j=1}^{K} \alpha_t(j) \beta_t(j)}
    \label{eq:gamma}
\end{equation}

\paragraph{M-Step.}
Re-estimate parameters using the posterior probabilities:
\begin{align}
    \mu_k^{\text{new}} &= \frac{\sum_{t=1}^{T} \gamma_t(k) \cdot O_t}{\sum_{t=1}^{T} \gamma_t(k)} \label{eq:mu_update} \\
    (\sigma_k^{\text{new}})^2 &= \frac{\sum_{t=1}^{T} \gamma_t(k) \cdot (O_t - \mu_k^{\text{new}})^2}{\sum_{t=1}^{T} \gamma_t(k)} \label{eq:sigma_update} \\
    T_{ij}^{\text{new}} &= \frac{\sum_{t=1}^{T-1} \xi_t(i,j)}{\sum_{t=1}^{T-1} \gamma_t(i)} \label{eq:T_update}
\end{align}
where $\xi_t(i,j) = \Prob(S_t = i, S_{t+1} = j \mid O_{1:T}, \mathcal{M})$.

\paragraph{Initialization.}
EM is sensitive to starting values.
We use a quantile-based initialization: observations are sorted and partitioned into $K$ equal-sized chunks, and each chunk's sample mean and standard deviation seed the corresponding state's emission parameters.
The transition matrix $\mathbf{T}$ and initial distribution $\pi$ are initialized uniformly.
This strategy prevents the common failure mode where all states collapse to the global mean.

\paragraph{Convergence.}
The algorithm terminates when the change in log-likelihood $\mathcal{L} = \text{logsumexp}(\log \alpha_T(\cdot))$ falls below $\text{tol} = 10^{-4}$ or after \texttt{max\_iter} iterations.

\subsection{Poisson Jump-Duration Mechanism}

To address the volatility-clustering limitation of standard HMMs \citep{ryden1998stylized, bulla2006stylized}, we augment the continuous HMM with the same Poisson jump-duration mechanism introduced by \citet{alswaidan2026hybrid}.
At each time step, with probability $\epsilon$, a jump event is triggered:
\begin{enumerate}
    \item A jump duration $D \sim \text{Poisson}(\lambda)$ is sampled.
    \item For each of the $D$ subsequent steps, the standard Markov transition is overridden.
    \item A coin flip with bias $p_{\text{neg}} = 0.52$ selects the target pool:
    \begin{itemize}
        \item \textbf{Crash states:} $\mathcal{S}^{-} = \{1, \ldots, N_{\text{tail}}\}$ (lowest-return regimes)
        \item \textbf{Boom states:} $\mathcal{S}^{+} = \{K - N_{\text{tail}} + 1, \ldots, K\}$ (highest-return regimes)
    \end{itemize}
    \item A state is drawn uniformly from the selected tail set.
\end{enumerate}

The coin flip occurs \emph{inside} the duration loop, so each jump step independently selects crash or boom.
This produces high volatility (magnitude) without directional trends, matching the empirical observation that large moves cluster regardless of sign \citep{mandelbrot1963variation, cont2001empirical}.

We denote the model without jumps as CHMM-NJ and the model with jumps as CHMM-WJ.

\subsection{Hyperparameter Optimization}

The jump parameters $\epsilon$ (tail-entry probability) and $\lambda$ (mean jump duration) are selected via a grid search that minimizes:
\begin{equation}
    J(\epsilon, \lambda) = \sum_{\tau=1}^{L} \left[\text{ACF}_{\text{obs}}^{|G|}(\tau) - \text{ACF}_{\text{sim}}^{|G|}(\tau)\right]^2 + w_K \left[K_{\text{obs}} - K_{\text{sim}}\right]^2
    \label{eq:objective}
\end{equation}
where $\text{ACF}^{|G|}(\tau)$ is the autocorrelation of absolute excess growth rates at lag $\tau$, $K_{\text{obs}}$ and $K_{\text{sim}}$ are the observed and simulated excess kurtosis, $L = 252$ trading days, and $w_K = 0.20$.
Simulated quantities are averaged over 200 independent paths per grid point.
The grid searches $\epsilon \in \{10^{-4}, 5 \times 10^{-4}, 10^{-3}, 5 \times 10^{-3}, 10^{-2}, 2.5 \times 10^{-2}\}$ and $\lambda \in \{10, 30, 55, 70, 85, 100, 130, 160\}$.

\subsection{Validation Metrics}

We evaluate synthetic data quality using five complementary metrics, following the framework of \citet{alswaidan2026hybrid} and \citet{stenger2024thinking}:

\begin{enumerate}
    \item \textbf{KS pass rate (\%):} Proportion of 1{,}000 simulated paths for which the two-sample Kolmogorov-Smirnov test \citep{kolmogorov1933, smirnov1948} fails to reject distributional equivalence at $\alpha = 0.05$.

    \item \textbf{Excess kurtosis:} Simulated excess kurtosis averaged across paths, compared to the observed value.
    Closer matching indicates better tail reproduction.

    \item \textbf{ACF-MAE:} Mean absolute error between the observed and simulated autocorrelation functions of absolute excess growth rates:
    \begin{equation}
        \text{ACF-MAE} = \frac{1}{L} \sum_{\tau=1}^{L} \left|\hat{\rho}^{\text{obs}}_{|G|}(\tau) - \hat{\rho}^{\text{sim}}_{|G|}(\tau)\right|
        \label{eq:acf_mae}
    \end{equation}
    Lower values indicate better reproduction of volatility clustering.

    \item \textbf{Wasserstein-1 distance ($W_1$):} Mean absolute difference of sorted empirical quantiles \citep{glasserman2003monte}.

    \item \textbf{Hellinger distance ($H$):} Histogram-based distributional overlap distance in $[0, 1]$; $H = 0$ indicates identical distributions.
\end{enumerate}
